% This is samplepaper.tex, a sample chapter demonstrating the
% LLNCS macro package for Springer Computer Science proceedings;
% Version 2.21 of 2022/01/12
%
\documentclass[runningheads]{llncs}
%
\usepackage[T1]{fontenc}
% T1 fonts will be used to generate the final print and online PDFs,
% so please use T1 fonts in your manuscript whenever possible.
% Other font encondings may result in incorrect characters.
%
\usepackage{graphicx}
% Used for displaying a sample figure. If possible, figure files should
% be included in EPS format.
%
% If you use the hyperref package, please uncomment the following two lines
% to display URLs in blue roman font according to Springer's eBook style:
%\usepackage{color}
%\renewcommand\UrlFont{\color{blue}\rmfamily}
%\urlstyle{rm}
%

\usepackage{amsmath}
\usepackage{amssymb}
\usepackage{cleveref}

\newcommand{\QQ}{\mathbb Q}
\newcommand{\ZZ}{\mathbb Z}


\begin{document}
%
\title{Implementing $p$-adic numbers in Macaulay2 using its foreign function interface and FLINT}


\titlerunning{Implementing $p$-adic numbers in M2 using FLINT}
%
\author{Douglas A. Torrance\inst{1}\orcidID{0000-0003-3999-4973}}
%
\authorrunning{D. Torrance}
% First names are abbreviated in the running head.
% If there are more than two authors, 'et al.' is used.
%
\institute{Piedmont University, Demorest GA 30535, USA
  \email{dtorrance@piedmont.edu}}

\maketitle              % typeset the header of the contribution
%
\begin{abstract}
  Macaulay2 is a computer algebra platform widely used by researchers in algebraic geometry and commutative algebra. Using the ForeignFunctions package, it is possible to make calls from Macaulay2 to dynamic libraries such as FLINT. We demonstrate this by introducing a new Macaulay2 package implementing p-adic numbers using FLINT via this interface. We discuss implementation details such as memory allocation, interaction with Macaulay2’s garbage collector, and object-oriented design decisions that mirror the existing implementations of the real and complex number fields in Macaulay2.
\keywords{Macaulay2  \and FLINT \and foreign function interface \and $p$-adic numbers}
\end{abstract}
%
%
%
\section{Introduction}

\section{A review of $p$-adic numbers}

We open with a crash course on $p$-adic numbers.  The reader is encouraged to reference a textbook such as \cite{gouvea} for more information.

Suppose $p$ is a prime number.  The field of \textit{$p$-adic numbers} $\QQ_p$ is a set whose elements may be expressed uniquely as formal Laurent series
\begin{equation}\label{p-adic expansion}
  \sum_{n=\nu}^\infty a_n p^n = a_\nu p^\nu + a_{\nu + 1}p^{\nu + 1} + \cdots,
\end{equation}
where $\nu\in\ZZ$ and $a_n\in\{0,\ldots,p-1\}$, together with the usual operations of addition and multiplication in base $p$, but noting that carrying may continue indefinitely.  An enlightening exercise is to verify that $-1 = \sum\limits_{n=0}^\infty(p - 1)\cdot p^n$.


For a given $x\in\QQ_p$, its \textit{$p$-adic valuation}, denoted $\nu_p(x)$, is the smallest nonzero exponent that appears in the expansion given in \Cref{p-adic expansion}, e.g., $\nu_3(1\cdot 3^{-2} + 2\cdot 3^{-1})=-2$.  In particular, $\nu_p(0) = \infty$.

These series are not convergent in general using the usual absolute value, but they do converge using the $p$-adic absolute value
\begin{equation*}
|x|_p = p^{-\nu_p(x)}.
\end{equation*}
In fact, $\QQ_p$ is the completion of $\QQ$ under $|\cdot|_p$.

The set of all $x\in\QQ_p$ with $\nu_p(x)\geq 0$ forms a ring, known as the \textit{$p$-adic integers} and denoted $\ZZ_p$.  Furthermore, the multiplicative subgroup $\ZZ_p^\times$ of the units of $\ZZ_p$ is precisely the $u\in\QQ_p$ for which $\nu_p(u) = 0$.  It follows that for any $x\in\QQ_p$ with $\nu_p(x)=\nu$,
\begin{equation}\label{in terms of unit}
  x = \sum_{n=\nu}^\infty a_n p^n = \left(\sum_{n=0}^\infty a_{n + \nu}p^n\right)p^\nu = up^\nu,
\end{equation}
where $u=\sum\limits_{n=0}^\infty a_{n + \nu}p^n\in\ZZ_p^\times$.

\begin{credits}
\subsubsection{\ackname} A bold run-in heading in small font size at the end of the paper is
used for general acknowledgments, for example: This study was funded
by X (grant number Y).

\subsubsection{\discintname}
It is now necessary to declare any competing interests or to specifically
state that the authors have no competing interests. Please place the
statement with a bold run-in heading in small font size beneath the
(optional) acknowledgments\footnote{If EquinOCS, our proceedings submission
system, is used, then the disclaimer can be provided directly in the system.},
for example: The authors have no competing interests to declare that are
relevant to the content of this article. Or: Author A has received research
grants from Company W. Author B has received a speaker honorarium from
Company X and owns stock in Company Y. Author C is a member of committee Z.
\end{credits}
%
% ---- Bibliography ----
%
% BibTeX users should specify bibliography style 'splncs04'.
% References will then be sorted and formatted in the correct style.
%
\bibliographystyle{splncs04}
\bibliography{macaulay2-padic}

\end{document}
